\begin{table}[t]
\centering
\caption{Auswahlverlust (\%) für Klassifikatoren und Abdeckung der Optima nach dem numerischen 5\%-Qualitätsfilter. Die 9 Gruppen werden gleich gewichtet. Die Abdeckung zählt Gruppen. Zufallswerte sind exakte Erwartungswerte. Budgets zählen nur Native-Tests der ausgeschlossenen Testgruppe. 151 verschiedene Trainingskonfigurationen über alle Folds, Generic-Messungen, Kompilierung und Qualitätsbewertung sind zusätzliche Kosten.}
\label{tab:supp-calibration-summary}
\begingroup\small\setlength{\tabcolsep}{3pt}
\par\smallskip\textbf{Auswahlverlust (\%)}\par\smallskip
\begin{tabular}{llll}
\toprule
Methode & k=1 & k=3 & k=5 \\
\midrule
Generic & 32.804 & 23.768 & 20.448 \\
Kalibriertes Generic & 28.319 & 23.146 & 15.431 \\
Gleichverteilte Zufallsauswahl & 32.290 & 15.406 & 9.104 \\
Grenzdatenmenge & 32.966 & 21.747 & 8.862 \\
\bottomrule
\end{tabular}
\par\smallskip\textbf{Optimum gefunden}\par\smallskip
\begin{tabular}{llll}
\toprule
Methode & k=1 & k=3 & k=5 \\
\midrule
Generic & 1 & 2 & 3 \\
Kalibriertes Generic & 1 & 3 & 4 \\
Gleichverteilte Zufallsauswahl & 0.603 & 1.808 & 3.014 \\
Grenzdatenmenge & 1 & 3 & 5 \\
\bottomrule
\end{tabular}
\endgroup
\end{table}
